\unnumberedchapter{چکیده}

بهره‌برداری از ریزشبکه یعنی در هر بازه زمانی معلوم شود بار از کجا تأمین شود، باتری شارژ گردد یا دشارژ، و با شبکه اصلی چقدر انرژی مبادله شود. این تصمیم‌ها باید با توازن توان و حدود ذخیره‌ساز سازگار باشند. اگر فقط یک شبکه عصبی روی داده آموزش ببیند، ممکن است خروجی از نظر عددی نزدیک داده باشد اما همان قیود را نقض کند. روش‌های بهینه‌سازی کلاسیک قیود را بهتر رعایت می‌کنند، ولی حل پیاپی آن‌ها همیشه به‌صرفه نیست.

در این پروژه از ایده شبکه‌های عصبی فیزیک‌آگاه استفاده شده است؛ یعنی باقیمانده توازن توان و پویایی وضعیت شارژ وارد تابع هزینه می‌شود. صورت‌بندی متغیرهای ریزشبکه و تبادل با شبکه از مدل‌های مدیریت انرژی متصل به شبکه گرفته شده است. لایه نظریه بازی برای تقسیم سود میان چند مشارکت‌کننده در آزمایش عددی وارد نشده و فقط در مرور منابع آمده است.

یک ریزشبکه ساده با تولید خورشیدی، باتری و بار محلی در نظر گرفته شد. شبکه پیش‌خور با پایتون پیاده گردید. داده از یک شبیه‌ساز ساعتی بیست‌ویک‌روزه با بذر ثابت به دست آمد. روی مجموعه آزمون، مدل فیزیک‌آگاه باقیمانده توازن و نقض حدود باتری را نسبت به مدل صرفاً داده‌محور کم کرد، هرچند خطای مربع میانگین کمی بیشتر شد. این مشاهده به همین داده مصنوعی محدود است.

\vspace{0.6cm}
\noindent\textbf{واژه‌های کلیدی:} \KeywordsFa
